\chapter{Computational algorithms}

\section{Algorithm 1: one-population channel decomposition}

\begin{lstlisting}
Input: SPD matrices G and H; channel dimension k.
1. Symmetrise numerical estimates: G <- (G + G^T)/2, H <- (H + H^T)/2.
2. Eigendecompose G = Q diag(g) Q^T.
3. Check all g_j > tolerance.
4. Form Ghalf = Q diag(sqrt(g)) Q^T.
5. Form S = Ghalf H Ghalf and symmetrise.
6. Eigendecompose S; order eigenvalues from largest to smallest.
7. Retain V = first k eigenvectors.
8. Map to phenotype space U = Ghalf V.
Output: eigenvalues, whitened basis V, phenotype basis U, spectral gap delta_k.
\end{lstlisting}

\section{Algorithm 2: invariant pairwise channel comparison}

\begin{lstlisting}
Input: phenotype bases U1 and U2; common SPD metric M.
1. A1 = U1^T M U1; A2 = U2^T M U2.
2. Q1 = U1 A1^{-1/2}; Q2 = U2 A2^{-1/2}.
3. Compute singular values sigma of Q1^T M Q2.
4. Clip numerical sigma to [0,1].
5. psi_j = arccos(sigma_j).
6. dC2 = sum sin(psi_j)^2.
7. IV = 1 - product cos(psi_j)^2.
Output: principal angles, chordal distance, volume deficit.
\end{lstlisting}

\section{Algorithm 3: operator-eigenspace susceptibility in a fixed representation}

\begin{lstlisting}
Input: ancestral S0; operator velocities A1 and A2; channel dimension k.
1. Eigendecompose S0 = [V W] diag(lambda) [V W]^T.
2. For every included b and excluded a:
       KS[a,b] = W[:,a]^T (A2-A1) V[:,b] /
                 (lambda[b]-lambda[k+a]).
3. chiS2 = sum KS[a,b]^2.
Output: KS and chiS2. This is the full coefficient when the whitening
        map is shared and fixed; otherwise continue with Algorithm 3b.
\end{lstlisting}

\section{Algorithm 3b: fully invariant susceptibility when \(\G\) evolves}

\begin{lstlisting}
Input: G0, H0, dG1, dH1, dG2, dH2, common metric M, dimension k.
1. R0 = sqrt(G0); S0 = R0 H0 R0.
2. Eigendecompose S0 and retain V, W, and eigenvalues lambda.
3. Solve R0 Xi + Xi R0 = dGi for X1 and X2.
4. For each population form
       dSi = Xi H0 R0 + R0 dHi R0 + R0 H0 Xi.
5. Compute KSi from dSi using the eigengap formula in Algorithm 3.
6. U0 = R0 V; B0 = sqrt(M) U0; Gamma0 = B0^T B0.
7. Q0 = B0 Gamma0^{-1/2}; construct an orthonormal complement Qperp.
8. KM = Qperp^T sqrt(M) [(X2-X1)V + R0 W(KS2-KS1)]
        Gamma0^{-1/2}.
9. chiM2 = sum KM[a,b]^2.
Output: KM and chiM2, where dC,M^2(t) = chiM2 t^2 + O(t^3).
\end{lstlisting}

\section{Algorithm 4: exact hybrid fitness}

\begin{lstlisting}
Input: hybrid mean mu, covariance P, environment optimum theta, curvature H.
1. d = mu - theta.
2. Expected log load:
       Llog = 0.5 d^T H d + 0.5 trace(H P).
3. For Gaussian arithmetic mean fitness:
       Hhalf = sqrt(H)
       C = Hhalf P Hhalf
       m = Hhalf d
       Wrel = det(I+C)^(-1/2) *
              exp[-0.5 m^T (I+C)^(-1) m].
Output: Llog and Wrel.
\end{lstlisting}

\section{Numerical stability}

\begin{itemize}[leftmargin=1.6em]
\item Use symmetric eigensolvers, not general eigensolvers, for \(\G\), \(\Hh\), and \(\Sop\).
\item Symmetrise estimates before decomposition to remove floating-point asymmetry.
\item Do not invert matrices explicitly when a linear solve is available.
\item Treat small or negative posterior eigenvalues carefully; they can indicate estimation noise or an invalid curvature model.
\item Compare subspaces rather than individual vectors near eigenvalue ties.
\item Propagate posterior uncertainty through every nonlinear operation.
\end{itemize}

\section{Independent numerical audit}

The companion script checks whitening, similarity of \(\G\Hh\) and \(\Sop\), generalized eigenvector residuals, common-metric coordinate invariance, projector and determinant identities, the fixed-representation perturbation coefficient, commuting perturbations, the full common-metric coefficient for evolving \(\G\) and \(\Hh\), coordinate invariance of that coefficient, both hybrid-fitness formulas, the general trace identity, the ridge theorem, the matrix-square-root derivative, the standing-variation response with migration, and the combined-barrier identity.

All nineteen checks pass under the fixed audit seed. The complete log is reproduced in \cref{app:numericalChecks}.
